% UpstreamNamespaces.tex -- namespaces of special functions in Mathlib and TauCeti,
% and the namespace policy of lean-DLMF.
% Compile with LuaLaTeX: latexmk -lualatex UpstreamNamespaces.tex
\documentclass[11pt,a4paper]{article}
\usepackage{leanreport}
\hypersetup{pdftitle={lean-DLMF: namespaces of special functions upstream, 2026-11-02}}

\begin{document}

\title{Namespaces of Special Functions\\in Mathlib and TauCeti}
\author{Bastiaan J. Braams}
\date{2026-11-02}
\maketitle

\begin{abstract}
  Many functions of the NIST Digital Library of Mathematical Functions (DLMF) already exist
  in Mathlib or TauCeti, and \texttt{lean-DLMF} either uses them directly or proves bridge
  theorems to them. This report surveys the namespaces in which Mathlib and TauCeti place
  special functions, polynomial families and number-theoretic functions. Mathlib follows a
  consistent rule: a declaration lives in the namespace of the object it is defined on or
  of the structure it belongs to. TauCeti mixes a \lean{TauCeti} prefix with bare Mathlib
  namespaces. Neither library has a \lean{SpecialFunctions} namespace. The report closes with
  the namespace policy of \texttt{lean-DLMF} and its consequences for bridges.
\end{abstract}

\section{Scope and method}

The survey covers Mathlib at revision \texttt{911003dd6d} and TauCeti at revision
\texttt{a049b4264d}, the versions pinned in \file{lake-manifest.json} of \texttt{lean-DLMF}
(Lean \texttt{v4.36.0-rc1}). Public definitions were collected from
\file{Mathlib/Analysis/SpecialFunctions}, \file{Mathlib/NumberTheory} (L-series, modular
forms, Bernoulli numbers, harmonic numbers, arithmetic functions, Dirichlet characters),
\file{Mathlib/RingTheory/Polynomial}, \file{Mathlib/Combinatorics/Enumerative}, and from
\file{TauCeti/Analysis/SpecialFunctions} and \file{TauCeti/NumberTheory}. The full name of
each definition was computed from the enclosing \texttt{namespace} blocks.

A preliminary remark: \file{Analysis/SpecialFunctions} is a directory, so it appears in
module names such as \module{Mathlib.Analysis.SpecialFunctions.Gamma.Basic}, but no
declaration in either library is in a namespace called \lean{SpecialFunctions}. Module paths
and declaration names are independent in Lean, and Mathlib chooses them independently.

\section{Mathlib}

\subsection{Analytic functions: the namespace of the domain}

A function of a real or complex variable lives in \lean{Real} or \lean{Complex}. When both
versions exist, each has its own namespace, and the real version is usually derived from the
complex one.

\begin{center}
\begin{tabularx}{\textwidth}{@{}l>{\RaggedRight\arraybackslash}X@{}}
  \toprule
  Namespace & Examples \\
  \midrule
  \lean{Complex} & \lean{Complex.Gamma}, \lean{Complex.digamma}, \lean{Complex.betaIntegral},
    \lean{Complex.GammaSeq}, \lean{Complex.Gammaℝ}, \lean{Complex.Gammaℂ},
    \lean{Complex.besselJ}, \lean{Complex.log}, \lean{Complex.cpow}, \lean{Complex.arctan},
    \lean{Complex.sinh}, \lean{Complex.regularizedHGFun} \\
  \lean{Real} & \lean{Real.Gamma}, \lean{Real.GammaSeq}, \lean{Real.doublingGamma},
    \lean{Real.arsinh}, \lean{Real.arcosh}, \lean{Real.artanh}, \lean{Real.logb},
    \lean{Real.arctan}, \lean{Real.eulerMascheroniConstant}, \lean{Real.binEntropy},
    \lean{Real.sigmoid} \\
  \lean{NNReal} & \lean{NNReal.agm} (arithmetic--geometric mean) \\
  \bottomrule
\end{tabularx}
\end{center}

\subsection{Functions at the root}

Several classical functions have no namespace at all. Most of them come from number theory,
where Mathlib treats the function as a mathematical object in its own right:
\lean{riemannZeta}, \lean{completedRiemannZeta}, \lean{hardyZ}, \lean{riemannZetaZeros},
\lean{jacobiTheta}, \lean{jacobiTheta₂}, \lean{ordinaryHypergeometric} (the Gauss series
$_2F_1$), \lean{bernoulli} (the Bernoulli numbers), \lean{bernoulliFun},
\lean{periodizedBernoulli}, \lean{harmonic}, \lean{catalan}, \lean{legendreSym},
\lean{jacobiSym} and \lean{gaussSum}.

\subsection{Functions attached to a structure or a theory}

When a function depends on a structure, it lives in the namespace of that structure:
\begin{itemize}
  \item \lean{PeriodPair.weierstrassP}: the Weierstrass $\wp$ function of a lattice given by a
    pair of periods;
  \item \lean{HurwitzZeta.hurwitzZeta} and \lean{HurwitzZeta.expZeta}, written as
    \lean{hurwitzZeta} after \texttt{open HurwitzZeta};
  \item \lean{ModularForm.eta} (Dedekind's $\eta$) and the \lean{EisensteinSeries} namespace;
  \item \lean{LSeries}, \lean{DirichletCharacter.LFunction}, \lean{ZMod.LFunction};
  \item \lean{Stirling.stirlingSeq}.
\end{itemize}

\subsection{Polynomial families}

Orthogonal and Appell polynomials are elements of a polynomial ring and live in
\lean{Polynomial}: \lean{Polynomial.Chebyshev.T} and \lean{Polynomial.Chebyshev.U},
\lean{Polynomial.hermite} (probabilists' normalization $\mathit{He}_n$) and
\lean{Polynomial.bernoulli}. The classical Jacobi, Gegenbauer and Laguerre families are not
in Mathlib.

\subsection{Number-theoretic and combinatorial functions}

These live in the namespace of the type they act on, or of their algebraic structure.

\begin{center}
\begin{tabularx}{\textwidth}{@{}l>{\RaggedRight\arraybackslash}X@{}}
  \toprule
  Namespace & Examples \\
  \midrule
  \lean{Nat} & \lean{Nat.totient}, \lean{Nat.primeCounting}, \lean{Nat.fib},
    \lean{Nat.centralBinom}, the Stirling numbers of both kinds \\
  \lean{ArithmeticFunction} & \lean{ArithmeticFunction.moebius},
    \lean{ArithmeticFunction.sigma}, \lean{ArithmeticFunction.vonMangoldt},
    \lean{ArithmeticFunction.liouville}, \lean{ArithmeticFunction.carmichael}; Dirichlet
    convolution is the ring multiplication \\
  \lean{Chebyshev} & \lean{Chebyshev.psi}, \lean{Chebyshev.theta} \\
  \lean{DirichletCharacter} & \lean{DirichletCharacter}, its conductor and primitivity \\
  \lean{Multiset} & \lean{Multiset.bell} \\
  \bottomrule
\end{tabularx}
\end{center}

\section{TauCeti}

TauCeti's own guideline is to put new material under a \lean{TauCeti} prefix, so that it
cannot collide with Mathlib. In practice three styles occur side by side.

\begin{center}
\begin{tabularx}{\textwidth}{@{}l>{\RaggedRight\arraybackslash}X@{}}
  \toprule
  Style & Examples \\
  \midrule
  \lean{TauCeti.} prefix & \lean{TauCeti.regularizedIncompleteBeta},
    \lean{TauCeti.inverseRegularizedIncompleteBeta}, \lean{TauCeti.lowerIncompleteGamma},
    \lean{TauCeti.regularizedGamma}, \lean{TauCeti.multivariateGamma},
    \lean{TauCeti.hermiteFunction}, \lean{TauCeti.chebyshevTHilbertBasis},
    \lean{TauCeti.hurwitzClassNumber}, \lean{TauCeti.BinaryQuadraticForm},
    \lean{TauCeti.Supernatural} \\
  \lean{TauCeti.Real.} prefix & \lean{TauCeti.Real.erf}, \lean{TauCeti.Real.erfc},
    \lean{TauCeti.Real.logIntegral} \\
  bare Mathlib namespace & \lean{Real.besselI}, \lean{Real.owensT}, \lean{Real.erfInv},
    \lean{DirichletCharacter.gaussSumOfPrimitiveRoot},
    \lean{DirichletCharacter.generalizedBernoulli}, \lean{QuadraticForm.finiteHasse},
    \lean{MonoidHom.cpowCharacter} \\
  \bottomrule
\end{tabularx}
\end{center}

The error function shows the mixture within one file: \lean{TauCeti.Real.erf} carries the
prefix, while its inverse \lean{Real.erfInv} does not. Declarations of the third style can
collide with a later Mathlib declaration of the same name, which is the situation the prefix
is meant to prevent. Most of TauCeti's number theory (binary quadratic forms, class numbers,
quadratic forms over number fields) uses the prefix; additions to existing Mathlib theories,
such as Dirichlet characters and quadratic forms, often do not.

\section{Consequences for \texttt{lean-DLMF}}

\subsection{The namespace policy}

The policy recorded in \file{AGENTS.md} on 2026-11-02 is the collision-free form of
TauCeti's guideline:
\begin{itemize}
  \item every declaration under \file{DLMF/} lives in \lean{DLMF.<Area>}, with real-variable
    forms in the \lean{Real} sub-namespace of the area, for example
    \lean{DLMF.Gamma.Real.digamma} beside Mathlib's \lean{Complex.digamma}; the only exception
    is a dot-notation lemma whose namespace is fixed by the head symbol of its conclusion;
  \item results in \file{DLMFMathlib/} extend Mathlib APIs, are not specific to DLMF, and use
    the Mathlib namespace of that API with Mathlib-style names, as candidates for
    upstreaming.
\end{itemize}
To bring the code into line, the Gamma area (112 declarations formerly in \lean{Real},
\lean{Complex} and \lean{MeasureTheory}) was moved into \lean{DLMF.Gamma} and
\lean{DLMF.Gamma.Real}, one root-level theorem was moved into \lean{DLMF.Hypergeometric}, and
two families of abbreviated root-level helpers in \file{DLMFMathlib/} were given the
descriptive namespaces \lean{StieltjesRemainder} and \lean{DoubleZeroSqrt}. All renames are
recorded in \file{reports/RenameMap20261101.tsv}.

\subsection{Functions that exist upstream}

When Mathlib or TauCeti already has a DLMF function, one of three patterns applies.
\begin{enumerate}
  \item The upstream definition matches DLMF in normalization, domain and branch. The project
    uses it directly, defines nothing locally, and records each DLMF item it covers with a
    checked pointer \texttt{dlmf\_upstream REF decl "relation: comment"}. Examples are
    \lean{Complex.besselJ}, \lean{riemannZeta}, \lean{PeriodPair.weierstrassP} and
    \lean{Complex.digamma}. There are about 1{,}000 such pointers.
  \item The upstream definition differs: real only, a series instead of the continuation,
    another normalization or another totalization. The project defines the DLMF function in
    \lean{DLMF.<Area>} and proves a bridge in the foundational layer. Examples are
    \lean{DLMF.Confluent.erf_ofReal} (to \lean{TauCeti.Real.erf}),
    \lean{DLMF.Bessel.besselI_ofReal} (to \lean{Real.besselI}), the Jacobi theta functions
    $\theta_j(z\,|\,\tau)$ defined from \lean{jacobiTheta₂}, and
    \lean{DLMF.Statistical.regularizedIncompleteBeta_ofReal} (to
    \lean{TauCeti.regularizedIncompleteBeta}). The logarithmic integral is bridged by
    $\operatorname{li}(x) = \operatorname{li}(2) + \operatorname{Li}(x)$ for $x > 1$, since
    \lean{TauCeti.Real.logIntegral} integrates from $2$ rather than taking DLMF's principal
    value from $0$.
  \item A thin wrapper \texttt{def X := upstream} is used only when a \texttt{@[dlmf]} tag
    needs a home that the upstream declaration cannot give, as for
    \lean{DLMF.Confluent.Real.invErf}.
\end{enumerate}

\subsection{Stability of bridges}

Bridges and pointers name the upstream declaration exactly as it is, so their statements
inherit the mixture described above: \lean{TauCeti.Real.erf}, \lean{Real.besselI} and
\lean{PeriodPair.weierstrassP} appear side by side. When TauCeti material moves into
Mathlib, its namespace is likely to change (\lean{TauCeti.Real.erf} would become
\lean{Real.erf}). The affected pointers and bridges then fail to compile, because
\texttt{dlmf\_upstream} checks that its target exists. This is the intended behaviour: the
build names every place that needs updating, and no pointer silently goes stale.

\end{document}
